% EXCERPT: original lines 1-70
\section{Separated stages and episode bounds}\label{sec:stages}

Fix \(\beta=1/2\). Suppose \eqref{eq:5} fails along integers
\(N\to\infty\) for some fixed \(D_0\), and let \(Q_N\) be the
law of the environment conditioned on \(a_N<N^{-\beta}\) along this
sequence. Then
\[
 {\rm KL}(Q_N\mid P)\le D_0\log N.
 \tag{34}\label{eq:34}
\]
We will divide the interval of heights from zero to \(N\) into episodes.
During an episode, a retained law of a fixed number of particles is repeatedly extended:
a successful extension increases its separation scale, and a failed
extension decreases that scale. An episode ends when the scale drops below
a fixed floor. The construction below makes failures very unlikely under
the fresh environment law \(P\). We then use the displayed entropy bound
to control their expected number under \(Q_N\).

\subsection{Separated extensions}

\begin{proposition}[Separated stages]\label{stage-prop}
There are fixed constants \(f,g,\chi,K>0\) and an integer \(k\ge2\)
with
\[
 \frac1g<\frac{\beta}{64},\qquad
 \frac{(1+f/g)D_0}{K}<\frac{\beta}{64},
\]
having the following property. For every sufficiently large fixed \(b\),
put \(B=kb\) and choose a sufficiently large fixed floor
\(s_{\rm fl}\). At every level-count scale \(s\ge s_{\rm fl}\), put
\[
 s_-=se^{-fb},\qquad s_+=se^{gb},\qquad
 H_e=\lfloor se^{-b}\rfloor,\qquad H=\lfloor se^{\chi b}\rfloor.
\]
For every probability \(\pi\) on \(k\)-tuples supported on
\(\{{\rm sep}\ge r_s\}\), a fresh product environment satisfies
\[
\begin{array}{ll}
 \Gamma_h^\pi\{{\rm sep}(\text{endpoints})\ge r_{s_-}\}\ge e^{-B}
       &(1\le h\le H),\\
 \Gamma_H^\pi\{{\rm sep}(\text{endpoints})\ge r_{s_+}\}\ge e^{-B}
\end{array}
\tag{35}\label{eq:35}
\]
with probability at least \(1-e^{-2Kb}\). Each test uses the full
kernel from the same stage base; a passing test does not restrict the
paths used by later tests.
\end{proposition}

\begin{proof}
The main construction gives several chances to create separated mass
before any prescribed short interval of target heights. A failed chance
reveals only layers below the next available chance, so the estimate can
be repeated there. We first describe these chances and their probability
bounds, and then choose the constants to cover all target heights.
Throughout these estimates \(A\ge1\) and the integer \(k\ge2\)
are provisional fixed parameters. Write
\[
 \Lambda=\log16,\qquad j_*=\frac1{8A},\qquad
 m=\lfloor j_*b\rfloor.
\]
We take \(b\) large enough that \(m\ge1\). All uses of the rapid-loss
bound \eqref{eq:32} will have cap \(B_*=2kb\); after fixing \(b\),
increasing \(s_{\rm fl}\) will ensure its separation hypotheses.

\paragraph{A nested schedule before one target interval.}
Cover \([H_e,H]\) by integer cells \([v,v+w_m]\), where
\[
 w_0=16^m\left\lfloor\frac{H_e}{4\cdot16^m}\right\rfloor,
 \qquad w_i=16^{-i}w_0\quad(0\le i\le m),

